% !TEX project = pdflatex
\documentclass[conference,a4paper]{IEEEtran}

\usepackage{packages}
\input{macros}

\begin{document}

\title{Big Data Analysis 1st Project (2025/26)\\
       \Large Track A: Spark SQL Optimization}

\author{
    \IEEEauthorblockN{André Miranda}
    \IEEEauthorblockA{MEI, PG60238}
    \and
    \IEEEauthorblockN{Gustavo Barros}
    \IEEEauthorblockA{MEI, PG61527}
    \and
    \IEEEauthorblockN{Miguel Carvalho}
    \IEEEauthorblockA{MEI, PG61153}
}

\maketitle

\begin{abstract}
The study demonstrates that targeted optimizations—reducing repeated dataset scans,
minimizing driver-side operations, and improving task parallelism—can yield
substantial performance gains in Spark-based HPC analyses. By combining internal
metrics from Spark event logs with system-level measurements from \texttt{dstat},
we quantify improvements in execution time, memory usage, and I/O throughput,
providing a comprehensive view of resource efficiency. The results highlight the
effectiveness of structured pipeline refactoring and offer practical guidance for
enhancing distributed data processing workflows in HPC environments.
\end{abstract}

\begin{IEEEkeywords}
Apache Spark, Spark SQL Optimization, HPC Job Analysis, Performance Benchmarking, Big Data Processing.
\end{IEEEkeywords}

\section{Introduction}
\label{sec:introduction}

High-Performance Computing (HPC) environments increasingly rely on distributed
data processing frameworks, such as Apache Spark, to handle large-scale datasets
efficiently. The combination of computational resources, parallel execution, and
intelligent data management allows complex analyses to be performed in reasonable
timeframes, but it also introduces challenges related to resource utilization,
I/O bottlenecks, and driver–executor coordination.

This report investigates the performance characteristics of a Spark-based Python
script, \texttt{statsEHPC\_v1\_init.py}, designed to analyze HPC job statistics
over multiple periods. The baseline implementation exhibited several performance
limitations, including repeated dataset scans, extensive driver-side operations,
and redundant shuffle activities. Understanding and mitigating these bottlenecks
is crucial for efficient data processing in HPC contexts, where system-level
contention and limited resources can exacerbate performance issues.

The methodology employed combines internal and external profiling tools. Execution
plans and Spark event logs were parsed using a custom \texttt{metrics\_parser.py}
script to extract task, stage, shuffle, and driver metrics. Simultaneously,
system-level utilization was measured with \texttt{dstat}, capturing CPU, memory,
disk I/O, and network usage during each run. Each optimization was validated
against the baseline output (\texttt{params\_baseline.tex}) to ensure correctness.

The report is organized as follows. Section~\ref{sec:provided_implementation}
describes the provided implementation, highlighting its key operations and
performance characteristics. Section~\ref{sec:methodology} details the benchmarking
approach and tools. Section~\ref{sec:optimizations} presents the optimizations,
their rationale, and the observed performance improvements. Finally,
Section~\ref{sec:analysis_results} discusses the results, identifying the most
effective strategies and analyzing the impact on resource utilization and execution
time.

\section{Provided Implementation}
\label{sec:provided_implementation}

The provided baseline implementation, \texttt{statsEHPC\_v1\_init.py}, defines
the reference pipeline for analysing HPC job accounting data and generating the
\LaTeX\ parameter file (\texttt{params.tex}) used in the final report.

The script initialises a Spark session and processes a set of monthly job log
files containing information about job execution, including job state,
partition, resource allocation, and runtime. Input files are discovered
dynamically and loaded sequentially, with each file being read using
\texttt{spark.read.csv()} and appended to a global DataFrame through iterative
\texttt{union} operations.

After ingestion, the dataset is enriched with derived attributes required for
the analysis. These include the extraction of the reporting period from the file
name, the classification of job outcomes into completed or failed states, and
the attribution of jobs to specific clusters and funding agencies.

Aggregated statistics are then computed for multiple reporting intervals (e.g.,
monthly, quarterly, and year-to-date). For each interval, the script repeatedly
applies \texttt{filter()} and \texttt{groupBy()} operations to compute metrics
such as the number of completed and failed jobs, total job counts, and total
execution time per cluster and funding category.

The results of these aggregations are transferred to the driver using multiple
\texttt{collect()} calls and stored in a dictionary. Finally, the dictionary is
written to \texttt{params.tex}, where each metric is exported as a \LaTeX\ macro.

While this implementation produces correct results, its structure introduces
several inefficiencies. In particular, sequential file loading, iterative
\texttt{union} construction, repeated dataset scans, and frequent
\texttt{collect()} operations limit the effectiveness of Spark’s distributed
execution model. These limitations motivate the optimisation strategies
presented in the following section.

\section{Methodology}
\label{sec:methodology}

The methodology followed in this project is based on a systematic cycle of
analysis, optimisation and evaluation of a Spark-based data processing pipeline
applied to HPC job accounting datasets.

The process begins with a detailed examination of the provided implementation,
focusing on its data flow, transformation pipeline and interaction with the
Spark execution model. Particular attention is given to how input data is
loaded, how derived attributes are computed, and how aggregations are executed
across different reporting periods. This initial inspection establishes a
baseline understanding of both the logical structure and the physical execution
behaviour of the application.

To support this analysis, Spark execution plans were collected using
\texttt{explain("extended")}, allowing inspection of both the logical and
physical plans generated by the Catalyst optimizer. In addition, Spark event
logging was enabled for all executions. These logs were processed using a
custom-built parser, which extracts key metrics such as the number of stages and
tasks, shuffle read/write volumes, and execution timing information. This
approach provides a fine-grained view of how transformations are translated into
distributed execution steps.

Resource utilization was measured externally using \texttt{dstat}, configured in
the benchmarking scripts to sample CPU, memory, network and disk I/O activity
during execution. CPU usage was collected using the \texttt{-tcmndry} flags,
capturing timestamp, CPU, memory, network, disk and page statistics. Memory
statistics capture usage and pressure over time, network throughput provides
insight into data transfer costs associated with shuffle operations, and disk
I/O metrics (reads/writes and throughput) help quantify the impact of data
spilling and interaction with the Parallel File System (PFS).

Each optimisation was implemented incrementally, following a controlled
experimental protocol. For every version of the pipeline, three independent runs
were executed under identical conditions, and average values were considered for
comparison. This ensures that observed improvements are consistent and not due
to transient system effects.

To guarantee correctness throughout the optimisation process, a validation step
was introduced. The output file \texttt{params.tex} generated by each version was
compared against a reference baseline (\texttt{params\_baseline.tex}) using the
\texttt{cmp} utility. This ensures that all optimisations preserve the exact
semantics and numerical results required by the assignment.

This methodology ensures that performance improvements are driven by measurable
evidence, combining execution plan analysis, event-level metrics, and system
resource monitoring, while strictly maintaining output correctness.

\section{Optimization Strategies}
\label{sec:optimizations}

The optimisation process was conducted incrementally, starting from the
baseline implementation and introducing targeted modifications aimed at
addressing specific inefficiencies in Spark’s execution model. Each version
focuses on a particular class of bottlenecks, including data ingestion,
redundant transformations, shuffle-heavy aggregations, and excessive
driver interaction.

\subsection{Baseline Bottleneck Analysis}
\label{sec:baseline_bottleneck_analysis}

Before introducing optimisations, the baseline implementation was analysed to
identify the main performance bottlenecks using execution time measurements,
Spark event logs, and inspection of the physical plan via
\texttt{explain("extended")}.

The following critical issues were identified:

\begin{itemize}
    \item \textbf{Iterative file loading and union building:} sequential file
    reads combined with repeated \texttt{union} operations create a deep
    lineage and unnecessary intermediate stages.

    \item \textbf{Repeated dataset scans:} independent
    \texttt{filter(...).groupBy(...)} patterns force Spark to rescan the same
    dataset multiple times.

    \item \textbf{Excessive shuffle operations:} multiple aggregations over the
    same keys trigger repeated wide transformations.

    \item \textbf{Frequent \texttt{collect()} calls:} repeated driver
    synchronisation introduces latency and limits scalability.

    \item \textbf{Schema inference overhead:} repeated schema inference during
    file loading adds unnecessary preprocessing cost.

    \item \textbf{I/O inefficiency on PFS:} multiple small writes degrade
    performance due to filesystem latency.
\end{itemize}

These bottlenecks guided the optimisation strategy described below.

\subsection{\texttt{v1} - Baseline}
\label{sec:v1}

The baseline implementation performs sequential ingestion, iterative unions,
repeated filtering and aggregation, and multiple driver-side collections.

\textbf{Why this is inefficient:}
This design underutilises Spark’s parallelism and forces repeated execution of
the same transformations due to lazy evaluation.

\textbf{Metrics (3 runs):}
\begin{itemize}
    \item Runtime: \texttt{\red{TODO}}
    \item \#Stages / \#Tasks: \texttt{\red{TODO}}
    \item Shuffle Read/Write: \texttt{\red{TODO}}
\end{itemize}

\textbf{Physical Plan Characteristics:}
\begin{itemize}
    \item Deep \texttt{Union} lineage
    \item Multiple \texttt{FileScan} and \texttt{Exchange} nodes
    \item Several independent jobs triggered by \texttt{collect()}
\end{itemize}

\subsection{\texttt{v2} - Parallel File Ingestion}
\label{sec:v2}

This optimisation replaces iterative file loading with a single
\texttt{spark.read.csv()} over all input files.

\textbf{Why:}
Iterative \texttt{union} builds a long lineage and introduces overhead in both
planning and execution. Spark is designed to parallelise file ingestion when
files are provided upfront.

\textbf{What is gained:}
\begin{itemize}
    \item Parallel file reading across executors
    \item Elimination of intermediate DataFrames
    \item Reduced DAG complexity and scheduling overhead
\end{itemize}

\textbf{Metrics (3 runs):}
\begin{itemize}
    \item Runtime: \texttt{\red{TODO}}
\end{itemize}

\textbf{Physical Plan Changes:}
\begin{itemize}
    \item Multiple \texttt{Union} nodes replaced by a single \texttt{FileScan}
    \item Flatter and more efficient lineage
\end{itemize}

\subsection{\texttt{v3} - Output I/O Consolidation}
\label{sec:v3}

This version buffers output and writes the parameter file once.

\textbf{Why:}
On HPC Parallel File Systems, small writes incur high latency due to metadata
operations and contention.

\textbf{What is gained:}
\begin{itemize}
    \item Reduced I/O overhead
    \item Lower filesystem contention
    \item Improved end-to-end runtime stability
\end{itemize}

\textbf{Metrics (3 runs):}
\begin{itemize}
    \item Runtime: \texttt{\red{TODO}}
\end{itemize}

\textbf{Physical Plan Impact:}
No direct change; improvement occurs outside Spark execution.

\subsection{\texttt{v4} - DataFrame Caching}
\label{sec:v4}

The main DataFrame is persisted using
\texttt{MEMORY\_AND\_DISK}.

\textbf{Why:}
Spark recomputes lineage for each action. When the same dataset is reused
multiple times, this leads to redundant computation and I/O.

\textbf{What is gained:}
\begin{itemize}
    \item Elimination of repeated dataset reconstruction
    \item Reduced CPU and I/O overhead
    \item More predictable execution times
\end{itemize}

\textbf{Metrics (3 runs):}
\begin{itemize}
    \item Runtime: \texttt{\red{TODO}}
\end{itemize}

\textbf{Physical Plan Changes:}
\begin{itemize}
    \item Introduction of \texttt{InMemoryTableScan}
    \item Removal of repeated upstream \texttt{FileScan}
\end{itemize}

\subsection{\texttt{v5} - Column Pruning}
\label{sec:v5}

Only required columns are selected after ingestion.

\textbf{Why:}
Processing unnecessary columns increases memory usage, serialization cost,
and shuffle volume.

\textbf{What is gained:}
\begin{itemize}
    \item Reduced I/O footprint
    \item Lower memory pressure
    \item More efficient shuffles
\end{itemize}

\textbf{Metrics (3 runs):}
\begin{itemize}
    \item Runtime: \texttt{\red{TODO}}
\end{itemize}

\textbf{Physical Plan Changes:}
\begin{itemize}
    \item Early \texttt{Project} operator
    \item Reduced data width across all stages
\end{itemize}

\subsection{\texttt{v6} - Filter Reuse}
\label{sec:v6}

Intermediate filtered DataFrames are reused.

\textbf{Why:}
Repeated filters force Spark to rescan and re-evaluate predicates on the same
dataset multiple times.

\textbf{What is gained:}
\begin{itemize}
    \item Reduced redundant scans
    \item Fewer predicate evaluations
    \item Improved logical plan compactness
\end{itemize}

\textbf{Metrics (3 runs):}
\begin{itemize}
    \item Runtime: \texttt{\red{TODO}}
\end{itemize}

\textbf{Physical Plan Changes:}
\begin{itemize}
    \item Fewer duplicated \texttt{Filter} nodes
    \item Reduced scan frequency
\end{itemize}

\subsection{\texttt{v7} - Aggregation Fusion}
\label{sec:v7}

Multiple aggregations are merged into a single
\texttt{groupBy().agg()}.

\textbf{Why:}
Each \texttt{groupBy} introduces a shuffle. Repeating it for each metric leads
to excessive data movement.

\textbf{What is gained:}
\begin{itemize}
    \item Significant reduction in shuffle operations
    \item Fewer Spark jobs and stages
    \item Improved network and disk efficiency
\end{itemize}

\textbf{Metrics (3 runs):}
\begin{itemize}
    \item Runtime: \texttt{\red{TODO}}
    \item Shuffle Read/Write: \texttt{\red{TODO}}
\end{itemize}

\textbf{Physical Plan Changes:}
\begin{itemize}
    \item Multiple \texttt{Exchange} nodes collapsed into one
    \item Single aggregation stage per dataset
\end{itemize}

\subsection{\texttt{v8} - Single-Pass Aggregation and Reduced \texttt{collect()}}
\label{sec:v8}

All metrics are computed in one aggregation and collected once.

\textbf{Why:}
Each \texttt{collect()} triggers a job and forces synchronisation between
executors and the driver.

\textbf{What is gained:}
\begin{itemize}
    \item Fewer driver–executor round-trips
    \item Reduced job scheduling overhead
    \item Better scalability in distributed execution
\end{itemize}

\textbf{Metrics (3 runs):}
\begin{itemize}
    \item Runtime: \texttt{\red{TODO}}
\end{itemize}

\textbf{Physical Plan Changes:}
\begin{itemize}
    \item Single aggregation and shuffle per iteration
    \item Elimination of multiple action-triggered jobs
\end{itemize}

\subsection{Summary}
\label{sec:optimizations_summary}

The optimisation process systematically eliminates the main inefficiencies of
the baseline implementation. The most impactful improvements arise from reducing
shuffle operations, avoiding repeated dataset scans, and minimising
driver-side interactions.

The final pipeline aligns with Spark’s execution model by favouring
single-pass transformations, reduced data movement, and efficient reuse of
intermediate results.

\section{Benchmark Policy}
\label{sec:benchmark_policy}

All experiments were executed on the ARM partition of the HPC cluster, with
a single node allocated per run. Each node provided 48 CPU cores and 16~GB of
RAM. Spark was configured to use 4 executors per node, each with 12 cores and
4~GB of executor memory. PySpark version 4.1.1 was used, and datasets were
read from the \texttt{/projects} directory on the Parallel File System (PFS),
where the benchmark output files were also written.

The benchmark procedure is automated via the \texttt{benchmark.sh} script,
located at the \texttt{jobs/} directory. For each optimization, the script
performs a dry run to warm up the JVM, then executes the target
\texttt{statsEHPC} Python script for three iterations, collecting system-level
resource utilization with \texttt{dstat} in CSV format (using the
\texttt{-tcmndry} flags to capture timestamp, CPU, memory, network, disk, and
page statistics). The generated Spark event logs are subsequently parsed using
\texttt{metrics\_parser.py} to extract application-level metrics, including
wall-clock time, stages, tasks, shuffle read/write volumes, driver time
estimate, and peak memory usage.

Peak CPU, memory, network, and disk I/O utilization are thus captured for each
iteration, enabling analysis of resource bottlenecks in addition to Spark's
internal metrics. The combination of event log analysis and system-level
monitoring ensures that both computational and I/O impacts are properly
quantified.

For consistency and correctness, the output generated by each run is validated
against the baseline results saved in \texttt{params\_baseline.tex} using
\texttt{cmp}, ensuring that optimizations do not alter the final computed
parameters.

\subsection*{Measured Metrics}
\label{sec:metrics}

The primary metrics collected for each benchmark run include:
\begin{itemize}
    \item \textbf{Wall time (s):} Application start to end.
    \item \textbf{Driver time (s):} Estimated driver-only time ($\text{wall} - \text{max-task-span}$).
    \item \textbf{Stages:} Number of stages submitted.
    \item \textbf{Tasks:} Total number of tasks executed across all stages.
    \item \textbf{Shuffle read/write (MB):} Volume of data read and written during
          shuffle operations.
    \item \textbf{Peak memory (MB):} Maximum memory consumption per node.
    \item \textbf{CPU, network, disk I/O:} System-level usage sampled by \texttt{dstat}.
\end{itemize}

All metrics are collected across three iterations, and the mean and standard
deviation are reported to capture variability. This procedure provides a
reliable basis for comparing the performance impact of each optimization.

\section{Aggregation of Results}
\label{sec:aggregation_results}

\red{TODO}

\section{Analysis of Results}
\label{sec:analysis_results}

The optimization process allowed a detailed examination of the performance characteristics
of the baseline implementation. Initial profiling revealed that the main bottlenecks
originated from repeated dataset scans, multiple driver-side operations such as
\texttt{collect()}, and excessive shuffle activity caused by repeated
\texttt{filter(...).groupBy(...)} sequences across periods.

By restructuring the computation pipeline, the following improvements were observed:

\begin{itemize}
    \item \textbf{Reduced I/O operations:} Prefiltering and avoiding repeated scans
          lowered the data read from the Parallel File System, reducing wall-clock
          time and network traffic.
    \item \textbf{Enhanced parallelism:} Partitioning and caching frequently accessed
          datasets improved task-level parallelism, leading to better CPU utilization
          and reduced driver wait times.
    \item \textbf{Minimized driver-side operations:} Shifting operations like
          aggregations and transformations to the executor side decreased memory
          pressure on the driver and eliminated repeated \texttt{collect()} calls.
\end{itemize}

The combination of these strategies led to a measurable decrease in execution time,
shuffle read/write volume, and peak memory usage, as reflected in the benchmark metrics
collected via \texttt{metrics\_parser.py} and \texttt{dstat}.

Among the optimizations, the most effective were those that reduced repeated scans
and shuffle operations, demonstrating that I/O-bound tasks were the dominant factor
in the original implementation's performance limitations. Optimizations targeting
executor-side computation and caching provided additional, albeit smaller, gains.

Finally, the validation of output using \texttt{params\_baseline.tex} confirmed that
all transformations preserved correctness, ensuring that performance improvements
did not compromise the final results.

\section{Conclusions}
\label{sec:conclusions}

This work analyzed the performance of a Spark-based HPC workflow and demonstrated
the impact of targeted optimizations on execution efficiency. The baseline
implementation suffered from repeated dataset scans, excessive driver-side
operations, and suboptimal parallelism, which were identified through a
combination of Spark event logs and system-level metrics collected with
\texttt{dstat}.

By applying pipeline restructuring, pre-filtering, and task-level optimizations,
we achieved significant reductions in wall-clock time and memory usage, while
maintaining result consistency, verified against the baseline outputs. The
benchmarking framework, including multiple iterations and automated validation
with \texttt{cmp} against \texttt{params\_baseline.tex}, ensured robust and
reproducible comparisons.

These results highlight the importance of understanding both the logical
execution plan and the physical resource utilization in HPC Spark applications.
Optimizations that reduce unnecessary I/O and improve task parallelism provide
the greatest gains. Future work may explore dynamic resource allocation, further
I/O tuning, and predictive scheduling strategies to anticipate and mitigate
potential bottlenecks in large-scale deployments.

\end{document}
